% Источники: 08_Протоколы_формальных_инспекций/DOC-13.5_Протокол_ФИ-АРХ-API.md
% Генератор: docs/build_pdf.py
% Шаблон печатного представления. Содержимое подставляет docs/build_pdf.py.
\documentclass[12pt,a4paper]{article}
\usepackage{iftex}
\ifPDFTeX
  \usepackage[utf8]{inputenc}
  \usepackage{cmap}
  \usepackage[T2A]{fontenc}
  % Tempora retains Times-like text and provides Cyrillic bold/italic shapes.
  \usepackage{tempora,newtxmath}
\else
  \usepackage{fontspec}
  \IfFontExistsTF{Times New Roman}
    {\setmainfont{Times New Roman}}
    {\setmainfont{Libertinus Serif}}
  \IfFontExistsTF{Arial}{\setsansfont{Arial}}{\setsansfont{Latin Modern Sans}}
  \IfFontExistsTF{Menlo}{\setmonofont{Menlo}[Scale=MatchLowercase]}
    {\setmonofont{Latin Modern Mono}[Scale=MatchLowercase]}
\fi
\usepackage[russian]{babel}
\usepackage{geometry}
\geometry{left=25mm,right=15mm,top=20mm,bottom=20mm}
\usepackage{setspace}
\onehalfspacing
\usepackage{indentfirst}
\setlength{\parindent}{1.25cm}
\setlength{\parskip}{2pt}
\setlength{\emergencystretch}{3em}
\usepackage{graphicx,booktabs,longtable,array,tabularx,calc,etoolbox}
\newcounter{none} % Pandoc 3.9 uses this counter for unnumbered longtables.
\usepackage{enumitem}
\setlist[itemize]{leftmargin=1.25cm,itemsep=2pt,topsep=4pt}
\setlist[enumerate]{leftmargin=1.25cm,itemsep=2pt,topsep=4pt}
\newcommand{\tightlist}{\setlength{\itemsep}{0pt}\setlength{\parskip}{0pt}}
\AtBeginEnvironment{longtable}{\small\singlespacing}
\setlength{\LTpre}{8pt}
\setlength{\LTpost}{8pt}
\setlength{\tabcolsep}{5pt}
\setlength{\LTleft}{0pt}
\setlength{\LTright}{0pt}
\usepackage{xurl}
\usepackage[unicode,colorlinks=true,linkcolor=black,urlcolor=blue]{hyperref}
\usepackage{bookmark}
\hypersetup{pdftitle={Протокол формальной инспекции архитектуры и API},pdfauthor={Анохин Г.А.},pdfsubject={Криптотренажёр mephi-12}}
\urlstyle{same}
\pagestyle{plain}
\setcounter{secnumdepth}{3}
\setcounter{tocdepth}{2}
\listfiles

\begin{document}
\hypersetup{pageanchor=false}
\begin{titlepage}
  \thispagestyle{empty}
  \centering
  {\large НИЯУ МИФИ\par}
  \vspace*{3.5cm}
  {\LARGE\bfseries Протокол формальной инспекции архитектуры и API\par}
  \vspace{0.5cm}
  {\large DOC-13.5\par}
  \vspace{0.6cm}
  {\large Платформа для изучения и практики курса по криптографии\par}
  {\large Криптотренажёр mephi-12\par}
  \vspace{0.7cm}
  {\large Версия 0.3\par}
  {\large Редакция от 25.11.2026\par}
  \vspace{1cm}
  \begin{tabularx}{\textwidth}{@{}p{3.4cm}Xp{4cm}@{}}
    \toprule
    Роль & ФИО & Подпись / Дата \\
    \midrule
    Разработал & Анохин Г.А. & \\
    \addlinespace[6pt]
    Согласовано & & \\
    \addlinespace[6pt]
    Утверждено & & \\
    \bottomrule
  \end{tabularx}
  \vfill
  {\large Москва, 2026\par}
\end{titlepage}
\hypersetup{pageanchor=true}

\section*{Лист регистрации изменений}
\addcontentsline{toc}{section}{Лист регистрации изменений}
\begingroup\small\singlespacing
\noindent
\begin{tabularx}{\textwidth}{@{}p{1.2cm}p{2cm}p{2.5cm}Xp{1.8cm}@{}}
\toprule
Версия & Дата & Автор & Содержание изменения & Утвердил \\
\midrule
0.3 & 25.11.2026 & Самбьюи А.Ж. & Сохранена исходная сверка от 20.11, добавлены материалы повторного разбора двух авторских исправлений; решение участников ожидается. & \\
\bottomrule
\end{tabularx}
\endgroup
\clearpage
\begingroup\small\singlespacing
\tableofcontents
\endgroup
\clearpage

\hypertarget{doc-135-протокол-формальной-инспекции-архитектуры-и-api}{}

Редакция 0.3 от 25.11.2026. Статус --- документарная сверка и материалы повторного разбора; решение участников инспекции ожидается. Исходная сверка редакции 0.2 от 20.11.2026 сохранена ниже как запись о рассмотренных тогда документах.

\section{Предмет проверки}\label{ux43fux440ux435ux434ux43cux435ux442-ux43fux440ux43eux432ux435ux440ux43aux438}

\hypertarget{предмет-проверки}{}

Проверены два документа версии 1.0 от 20.11.2026. В DOC-07 «Архитектура ПО» указано, какой компонент отвечает за действие. В DOC-08 «Справочник API платформы» описан программный интерфейс (API) --- какие запросы принимает сервер и какой результат получает пользователь. При сверке прослежены действия преподавателя и студента от публикации контрольной до чтения результатов. Для спорных случаев приведены данные, на которых два описания дают разные результаты.

{\def\LTcaptype{none} % do not increment counter
\begin{longtable}[]{@{}
  >{\raggedright\arraybackslash}p{(\linewidth - 2\tabcolsep) * \real{0.3200}}
  >{\raggedright\arraybackslash}p{(\linewidth - 2\tabcolsep) * \real{0.6800}}@{}}
\toprule\noalign{}
\begin{minipage}[b]{\linewidth}\raggedright
Проверенная редакция
\end{minipage} & \begin{minipage}[b]{\linewidth}\raggedright
Рассмотренные разделы
\end{minipage} \\
\midrule\noalign{}
\endhead
\bottomrule\noalign{}
\endlastfoot
{\href{../09\_Архитектура\_и\_API/DOC-07\_Архитектура.pdf\#doc-07-архитектура-по}{DOC-07, версия 1.0}} & Раздел 2 «Слои и компоненты», разделы 4--6 о записи, сроках и восстановлении \\
{\href{../09\_Архитектура\_и\_API/DOC-08\_Справочник\_API.pdf\#doc-08-справочник-api-платформы}{DOC-08, версия 1.0}}, включая {\href{../09\_Архитектура\_и\_API/DOC-08\_Справочник\_API.pdf\#doc-08-справочник-api-сессии-задач}{«Сессии задач»}} & Публикация, расписание, результаты, учётные записи и группы. В части о сессиях --- создание, отправка решения, чтение и ошибки \\
\end{longtable}
}

Для сравнения использованы карточки DOC-10 «План верификации ПО». Карточка с обозначением TC задаёт условия проверки, действия и ожидаемый результат. Ниже зафиксированы выводы по текстам документов.

\section{Результаты сверки}\label{ux440ux435ux437ux443ux43bux44cux442ux430ux442ux44b-ux441ux432ux435ux440ux43aux438}

\hypertarget{результаты-сверки}{}

\textbf{Группы и доступ к результатам.} Номер учебной группы студента сервер получает при регистрации или привязке подтверждённой учётной записи МИФИ. Группы преподавателю назначает администратор. Если администратор снимает назначение после публикации, преподаватель при следующем запросе результатов своей работы для этой группы получает HTTP 404. Правило доступа задано в DOC-07, раздел 2. Код ответа указан в DOC-08, «Результаты контрольного мероприятия». Подтверждение группы и отзыв назначения разобраны в карточках {\href{../10\_Верификация\_и\_тестирование/DOC-10\_План\_верификации.pdf\#tc-auth-004}{TC-AUTH-004}} и {\href{../10\_Верификация\_и\_тестирование/DOC-10\_План\_верификации.pdf\#tc-auth-005}{TC-AUTH-005}}.

\textbf{Закрытие работы.} В карточке TC-SES-017 студент начинает работу в 10:55, а в 11:00 она закрывается. Сервер отклоняет новый ответ, полученный ровно в 11:00, с кодом 410 \texttt{SESSION\_\allowbreak{}EXPIRED}. Если полный допустимый ответ получен в 10:59:59.900, вердикт по нему можно записать в 11:00:02. Ровно в 11:00:59.900 срок записи уже истёк. DOC-07, раздел 5, и DOC-08, «Сессии задач», раздел 2, дают тот же результат. Связь с системными требованиями SR-SES-011 и SR-SES-012 указана в {\href{../10\_Верификация\_и\_тестирование/DOC-10\_План\_верификации.pdf\#соответствие-карточек-требованиям}{таблице соответствия DOC-10}}.

\textbf{Возврат после обрыва связи.} Если студент потерял ответ на запрос старта, в расписании он находит идентификатор своей сессии и открывает её запросом чтения. После завершения сессии новый запрос старта при открытой контрольной и разрешённых повторах создаёт отдельное прохождение. Такое различие явно описано в DOC-08, «Расписание студента» и «Сессии задач», раздел 1. Карточка TC-CFG-005 проверяет условия повторного старта.

\section{Замечания к разбору}\label{ux437ux430ux43cux435ux447ux430ux43dux438ux44f-ux43a-ux440ux430ux437ux431ux43eux440ux443}

\hypertarget{замечания-к-разбору}{}

\subsection{Состояние просроченной задачи в ведомости преподавателя}\label{1-ux441ux43eux441ux442ux43eux44fux43dux438ux435-ux43fux440ux43eux441ux440ux43eux447ux435ux43dux43dux43eux439-ux437ux430ux434ux430ux447ux438-ux432-ux432ux435ux434ux43eux43cux43eux441ux442ux438-ux43fux440ux435ux43fux43eux434ux430ux432ux430ux442ux435ux43bux44f}

\hypertarget{1-состояние-просроченной-задачи-в-ведомости-преподавателя}{}

В DOC-07, раздел 5, и DOC-08, «Сессии задач», раздел 3, студент при чтении своей сессии сразу видит просрочку. В DOC-08, «Результаты контрольного мероприятия», описан состав ведомости преподавателя, но не указано, на какой момент определять состояние задач.

Например, работа закрылась в 10:30, студент ничего не отправил, фоновая проверка запланирована на 10:30:30. В 10:30:01 студент уже увидит у задачи \texttt{FAILED} (срок истёк) и \texttt{result\ =\ null} (вердикта нет). В ведомости преподавателя по тексту допустимо показать прежнее \texttt{NEW} --- задача ещё открыта. От этого зависит, какую работу преподаватель сочтёт незавершённой.

Предложение редактору DOC-08 --- установить в описании ведомости тот же момент определения состояния, что и при чтении студентом. Для повторной сверки нужны два случая между закрытием и ближайшей фоновой проверкой. При отсутствии ответа обе стороны видят просрочку. Если своевременный ответ ещё обрабатывается в пределах отведённых 60 секунд, обе стороны видят открытую задачу до записи вердикта либо истечения срока обработки.

\subsection{Условия проверки публикации пустой контрольной}\label{2-ux443ux441ux43bux43eux432ux438ux44f-ux43fux440ux43eux432ux435ux440ux43aux438-ux43fux443ux431ux43bux438ux43aux430ux446ux438ux438-ux43fux443ux441ux442ux43eux439-ux43aux43eux43dux442ux440ux43eux43bux44cux43dux43eux439}

\hypertarget{2-условия-проверки-публикации-пустой-контрольной}{}

В {\href{../10\_Верификация\_и\_тестирование/DOC-10\_План\_верификации.pdf\#tc-cfg-008}{карточке TC-CFG-008}} случай «ноль и одно задание» предлагает передать группу и время открытия пустой работы. Ожидается код 409 \texttt{ASSESSMENT\_\allowbreak{}EMPTY}. По DOC-08, «Назначение группе и публикация», при пропуске времени закрытия сервер вернёт 400 ещё до проверки состава. Далее карточка ожидает успешную публикацию после добавления задания, хотя в описанных действиях преподаватель ещё не выбрал правило повторного прохождения.

Предложение владельцу DOC-10 --- во всех запросах этого случая указывать допустимые \texttt{groupId}, \texttt{opensAt} и \texttt{closesAt}. Для пустого состава ожидается 409 \texttt{ASSESSMENT\_\allowbreak{}EMPTY}. После добавления задания при невыбранном правиле ожидается 409 \texttt{ASSESSMENT\_\allowbreak{}RETAKE\_\allowbreak{}POLICY\_\allowbreak{}REQUIRED}. После явного выбора, например \texttt{repeat\allowbreak{}Allowed =\allowbreak{} false}, ожидается 200 и карточка появляется в расписании назначенной группы. Для двух параллельных запросов в конце той же карточки выбор повтора также должен быть задан заранее.

\section{Состав инспекции и решение}\label{ux441ux43eux441ux442ux430ux432-ux438ux43dux441ux43fux435ux43aux446ux438ux438-ux438-ux440ux435ux448ux435ux43dux438ux435}

\hypertarget{состав-инспекции-и-решение}{}

По графику процедуры DOC-04 заседание планировалось на 16--17.11.2026, выпуск протокола --- на 17.11.2026. Рассмотренные редакции DOC-07 и DOC-08 датированы 20.11.2026. В графике владельцами документов указаны Георгиев Д.Д. и Минасов М.Р., ведущим --- Портнов Л.М., секретарём --- Анохин Г.А., рецензентами --- Кокорев В.Д. и Самбьюи А.Ж. {\href{../../06\_Процедура\_формальной\_инспекции/Процедура\_формальной\_инспекции.pdf}{Процедура и график инспекций, DOC-04}}.

При назначении секретаря на эти редакции ведущему нужно учесть авторство. Анохин Г.А. указан в DOC-07 как подготовивший редакцию и в DOC-08 как автор раздела о сессиях. DOC-04, раздел «Роли и ответственности», поручает протокол заседания участнику, который не писал проверяемый документ. Для совместного рассмотрения этих редакций требуется другой секретарь.

На 20.11.2026 документарная сверка дала два замечания. Фактическая дата заседания, явка, листы рецензентов и подписи участников не подтверждены. Решение по обоим замечаниям и итог приёмки остаются за участниками инспекции. Назначенный секретарь внесёт в этот протокол решения, сроки исправления и номера принятых замечаний в общем реестре (DEF). Число принятых замечаний и показатели заседания можно определить после его оформления.

{\def\LTcaptype{none} % do not increment counter
\begin{longtable}[]{@{}
  >{\raggedright\arraybackslash}p{(\linewidth - 2\tabcolsep) * \real{0.3200}}
  >{\raggedright\arraybackslash}p{(\linewidth - 2\tabcolsep) * \real{0.6800}}@{}}
\toprule\noalign{}
\begin{minipage}[b]{\linewidth}\raggedright
Запись по итогам заседания
\end{minipage} & \begin{minipage}[b]{\linewidth}\raggedright
Сведения на 20.11.2026
\end{minipage} \\
\midrule\noalign{}
\endhead
\bottomrule\noalign{}
\endlastfoot
Дата, время, место проведения и фактический состав & Ожидают подтверждения ведущим \\
Решение по замечаниям 1 и 2, номера DEF, ответственные и сроки & Ожидают разбора на заседании \\
Решение по DOC-07 и DOC-08, дата повторной инспекции при необходимости & Решение ведущего в материалах сверки отсутствует \\
Листы рецензентов, подписи и даты & В материалах сверки отсутствуют \\
\end{longtable}
}

\section{Дополнение от 25.11.2026: материалы повторного разбора}\label{ux434ux43eux43fux43eux43bux43dux435ux43dux438ux435-ux43eux442-25112026-ux43cux430ux442ux435ux440ux438ux430ux43bux44b-ux43fux43eux432ux442ux43eux440ux43dux43eux433ux43e-ux440ux430ux437ux431ux43eux440ux430}

\hypertarget{дополнение-от-25112026-материалы-повторного-разбора}{}

Подготовил Самбьюи А.Ж. в объёме рецензента задачи 3.6 и автора собственных исправлений. Исторические выводы о версиях от 20.11 и сведения об отсутствии заседания выше сохраняются. Текущие объекты повторного рассмотрения --- {\href{../09\_Архитектура\_и\_API/DOC-08\_Справочник\_API.pdf\#doc-08-справочник-api-платформы}{DOC-08 v1.2}} и {\href{../10\_Верификация\_и\_тестирование/DOC-10\_План\_верификации.pdf\#doc-10-план-верификации-программного-обеспечения}{DOC-10 v1.1}}; архитектурные основания остаются в DOC-07.

{\def\LTcaptype{none} % do not increment counter
\begin{longtable}[]{@{}
  >{\raggedright\arraybackslash}p{(\linewidth - 4\tabcolsep) * \real{0.2359}}
  >{\raggedright\arraybackslash}p{(\linewidth - 4\tabcolsep) * \real{0.3768}}
  >{\raggedright\arraybackslash}p{(\linewidth - 4\tabcolsep) * \real{0.3872}}@{}}
\toprule\noalign{}
\begin{minipage}[b]{\linewidth}\raggedright
Предложение исходной сверки
\end{minipage} & \begin{minipage}[b]{\linewidth}\raggedright
Авторское изменение
\end{minipage} & \begin{minipage}[b]{\linewidth}\raggedright
Повторная независимая проверка
\end{minipage} \\
\midrule\noalign{}
\endhead
\bottomrule\noalign{}
\endlastfoot
1. Просрочка в ведомости преподавателя & DOC-08 v1.1 задаёт состояние на момент чтения по правилам своей сессии; текст сохраняется в v1.2 & Сопоставить три представления в D, после D и в L; сохранённый вердикт не подменяется просрочкой \\
2. Предусловия TC-CFG-008 & DOC-10 v1.1 задаёт группу, opensAt, closesAt, затем пустой состав, отсутствие выбора повтора и явное false & Убедиться, что запросы различаются только рассматриваемым условием; проверить порядок двух конкурирующих записей \\
\end{longtable}
}

Оба пункта имеют статус «текст исправлен автором; принятие участниками ожидается». Номера DEF, закрытие замечаний и решение ФИ этим дополнением не назначаются.

Самбьюи подготовил собственную часть API и изменения кейсов; он не выдаёт их самооценку за независимую рецензию. Ведущему следует поручить проверку этих исправлений другому участнику. Для чужих архитектурных и API-разделов Самбьюи может подготовить лист замечаний с точным местом и источником; подтверждённых новых замечаний к ним в данном дополнении не заявляется. Конфликт авторства секретаря исходных редакций, отмеченный выше, также требует разрешения ведущим.

Для повторного разбора подготовлены проверяемые случаи: задача без ответа сразу после закрытия до фонового прохода; своевременная обработка между D и L; точный L; сохранённый до L вердикт; таймаут L до D; доступ после отзыва группы. Ожидаемые реакции приведены в DOC-08 и TC-CFG-002. Это материалы проверки документации; испытания программного обеспечения не проводились.


\end{document}
